\chapter{Anhang}

\section{Oszilloskopaufnahmen Szintillationsdetektor-Signale}

\jafps{TEK0002}{Oszilloskopaufnahme eines weiteren Anodendraht-Signals (in orange). Anpassungsparameter an die Anpassungsfunktion (in blau): $a=\SI{-4.668(61)e+05}{\per\s}, b=\SI{3.85(45)}{\milli\V}, c=\SI{-0.0(3.6)e+04}{\V}, t_0=\SI{-0.00(32)}{\s}$, als Dauer (grün hinterlegt) $d=\SI{1.8720(28)}{\micro\s}$.}{0.5}
\jafps{TEK0003}{Oszilloskopaufnahme eines weiteren Anodendraht-Signals (in orange). Anpassungsparameter an die Anpassungsfunktion (in blau): $a=\SI{-5.76(13)e5}{\per\s}, b=\SI{1.26(52)}{\milli\V}, c=\SI{-0.0(4.8)e4}{\V}, t_0=\SI{-0.00(40)}{\s}$, als Dauer (grün hinterlegt) $d=\SI{1.3920(28)}{\micro\s}$.}{0.5}

\section{Strommessung}

\jafps{strom_ohne}{Darstellung des Stroms durch die Driftkammer $I_D$ ohne Quelle. Man erkennt, dass sich der Strom im Rahmen der Messunsicherheit nur sehr geringfügig mit $U_D$ verändert.}{0.5}


\begin{table}[h]
\centering
\begin{tabular}{|c|c|} \hline
Drahtnummer & Winkel / Grad \\\hline
$0$ & $\num{-60.5(3.4)}$ \\
$2$ & $\num{-58.2(3.9)}$ \\
$4$ & $\num{-55.5(4.4)}$ \\
$6$ & $\num{-52.4(5.1)}$ \\
$8$ & $\num{-48.8(5.9)}$ \\
$10$ & $\num{-44.5(6.9)}$ \\
$12$ & $\num{-39.6(8.0)}$ \\
$14$ & $\num{-33.8(9.2)}$ \\
$16$ & $\num{-27(11)}$ \\
$18$ & $\num{-20(12)}$ \\
$20$ & $\num{-11(13)}$ \\
$22$ & $\num{-2(13)}$ \\
$24$ & $\num{7(13)}$ \\
$26$ & $\num{15(12)}$ \\
$28$ & $\num{23(11)}$ \\
$30$ & $\num{30.6(9.9)}$ \\
$32$ & $\num{36.8(8.6)}$ \\
$34$ & $\num{42.1(7.4)}$ \\
$36$ & $\num{46.7(6.4)}$ \\
$38$ & $\num{50.7(5.5)}$ \\
$40$ & $\num{54.0(4.7)}$ \\
$42$ & $\num{56.9(4.1)}$ \\
$44$ & $\num{59.4(3.6)}$ \\
$46$ & $\num{61.6(3.2)}$ \\
$48$ & $\num{63.5(2.8)}$ \\
\hline
\end{tabular}
\caption{Winkel der Bins mit Unsicherheiten.}
\label{tab:angle_bins}
\end{table}
